
\documentclass{article}
\usepackage[utf8]{inputenc}
\usepackage{amsmath}
\usepackage{fullpage}
\usepackage{xcolor}
\usepackage{float}
\usepackage{graphicx}
\usepackage{soul}
\usepackage{ulem}
\usepackage{amsfonts}

\newcommand{\mat}[1]{\mathbf{#1}}
\newcommand{\ten}[1]{\mathbf{\mathcal{#1}}}
\newcommand{\response}[1]{\textbf{Response:} \textcolor{blue}{\textit{#1}}\\}
\newcommand{\GB}[1]{\textcolor{red}{GB: #1}}

\newcommand\numberthis{\addtocounter{equation}{1}\tag{\theequation}}
\title{Responses to Reviewer's Comments}
% \author{bhishamdevverma }
\date{}

\begin{document}

\maketitle
% \vspace{-1.5cm}

We thank the reviewers for dedicating their time and effort to reviewing the manuscript thoroughly and providing insightful comments. The manuscript has significantly benefited from them. All of the changes to the manuscript are colored in blue.  We provide responses to individual comments below.\\


\noindent \underline{\textbf{Referee \#1}}\\


\noindent \ul{\textit{\textbf{Suggestion-1:}}}  In several places, when objects are multiplied/applied, they have compatible sizes but are of the wrong shapes.
For example, in (1.3), the matrices $\mat{Q}_1$, $\mat{Q}_2$, and $\mat{Q}_3$ are $n \times r$, so the operator $(\mat{Q}_1 \otimes \mat{Q}_2 \otimes \mat{Q}_3)$ should act on a object with "column" shape $r \times r \times r$.
The matrix $\hat{\mat{Q}}_{12}$ is $r^2 \times r$ and $\mat{I}_r$ is $r \times r$, so the object $(\hat{\mat{Q}}_{12} \otimes I_r)$ has "column" shape $r^2 \times r$.
So, in principle, the product $(\mat{Q}_1 \otimes \mat{Q}_2 \otimes \mat{Q}_3)(\hat{\mat{Q}}_{12} \otimes \mat{I}_r)$ is incorrect.
To make all the equations formally correct would require introducing explicit reshaping operators, which would then make all the equations much harder to read.
I suggest to instead explain this size/shape issue and then say that objects implicitly reshape themselves as needed.\\

\response{We regret the confusion, we mean to use $\otimes$ to denote the matrix Kronecker product, so that all of these expressions are matrices, and their dimensions are consistent for matrix multiplication.  To clarify this in the paper, we make explicit that the notation is for matrix Kronecker products both in the introduction, where we give a teaser for the algebra, and in the Background section, where we define things formally.  We also replaced a separate use of $\otimes$ that did denote a tensor product to avoid confusion.}

\noindent \ul{\textit{\textbf{Suggestion-2:}}} p13 line 290 gives [4,10] as references for the sine of sums tensor identity (4.1). Add a citation to the paper with the proof http://dx.doi.org/10.1007/BF02985795\\

\response{Done.}


\noindent \underline{\textbf{Referee \#2:}} \\

\noindent\underline{General comments: (comments on improving the manuscript in general)}\\

\noindent \ul{\textit{\textbf{GC-1:}}} The references could be filled out a bit more, e.g., adding references like Khatri and Rao's paper from 1968, etc.\\

\response{Added references to original Khatri-Rao paper and CP papers.}

\noindent \ul{\textit{\textbf{GC-2:}}} I think the paper would be made stronger with an example of how the presented algorithms perform on an application (e.g., a selected problem from the list of applications presented in the introduction). This would not be a minor revision, and is not required, but I think the paper could benefit from having such an example.\\

\response{We agree that demonstration of the presented algorithms on another application would make the paper stronger.  We have elected not to pursue another example because we expect the performance behavior to be qualitatively similar to what is observed in the existing examples.  The accuracy behavior is harder to predict, as it depends on the condition numbers coming from the problem.  We assume for most applications, the conditioning of the problem will be okay for the NE-based approach, at least in double precision.  But because the QR-based approach is likely to be just as fast or faster, we still recommend its use in general.}

\noindent \ul{\textit{\textbf{GC-3:}}} Background or tutorial references on the complexity analysis of some of the operations more specific to tensor decompositions (e.g., MTTKRP, Multi-TTM, use of the arithmetic average and geometric mean) would be helpful, as they aren't obvious to non-experts in the field.\\

\response{Using arithmetic average and geometric mean is strictly for notational convenience and presentation.  The complexity of these computations typically depends on either the product of the dimensions or the sum, so using the average and geometric mean avoids the cumbersome sum and product notation.  It also makes the simplification of the expressions for cubical tensors ($n_1=\cdots=n_d$) obvious.  We have added references to justify the costs of TTM, MTTKRP, and Multi-TTM throughout Section 3 and point readers to tutorials on those computations.}

\noindent \underline{Specific comments:}\\

\noindent \ul{\textit{\textbf{SC-1:}}} I would state Property 1, with reference, before (1.1)-(1.3).\\

\response{We have added a forward pointer to Property 1 in the introduction before (1.1)-(1.3).  We would like to avoid overloading the introduction with mathematical results, so we left Property 1 in the background section, but we did want to present the key mathematical idea as a teaser.}

\noindent \ul{\textit{\textbf{SC-2:}}} Memoization is mentioned on page two without reference, but later is referenced (with two different references).\\

\response{This is intentional, memoization is a generic technique.   We have rephrased the sentence in the introduction to try and clarify this.  The way we use it as described on page 2 is novel because it is in the context of our algorithm, so we don't think a reference is appropriate.  We use a reference in Section 3.1.1 because that paper applied memoization in exactly that context (CP-ALS of a dense tensor).  We use a different reference in Section 3.2.1 because that paper addressed how to use memoization in computing a sequence of Multi-TTM computations, though in a slightly different context.}

\noindent \ul{\textit{\textbf{SC-3:}}} There was a bit of a disconnect between complexities presented in Table 2 and subsection 3.1.1 for me that was resolved by understanding that the table presents per-iteration complexity whereas the discussion in 3.1.1 centers around sub-iterations. A more clear explanation of what is an iteration vs. sub-iteration would probably help some readers. This could be done in section 2.3.\\

\response{Done.}

\noindent \ul{\textit{\textbf{SC-4:}}} Computing ``$\mat{H}_k$ in line 8" should specify that you're talking about line 8 of Algorithm 3.1.\\

\response{Done.}

\noindent \ul{\textit{\textbf{SC-5:}}} The algorithms contain while loops with the condition ``not converged". This is vague; in what sense have the algorithms converged?\\

\response{We added a short discussion of convergence criteria in the beginning of Section 3. Basically, convergence criteria are orthogonal to the per-iteration complexities; changing the convergence criteria does not affect the qualitative comparisons among algorithms.  In our experiments we run to a maximum number of iterations for simplicity.}

\noindent \ul{\textit{\textbf{SC-6:}}} In algorithm 3.3, why is the expression for script $\ten{Y}$ used on line 8 when $\mat{Y}_{(k)}$ is used on line 10. To me, using the above expression for $\mat{Y}_{(k)}$ would be more clear. Is this because the notation is more convenient for later algorithms?\\

\response{We switch between tensor notation and matrix notation in the algorithm to make the computation clear.  $\mat{Y}_{(k)}$ is just the $k$th modal unfolding of $\ten{Y}$.  We can write line 8 in matrix notation (it includes the Kronecker products of the $\mat{Q}$'s and appears on page 8), but we want to emphasize that the efficient computation is a Multi-TTM.  We have added a comment in the background section to clarify that $\mat{Y}_{(k)}$ is the notation for unfolding.}

\noindent \ul{\textit{\textbf{SC-7:}}} On page 9, reference [7] is used for memoization, whereas [11] was used in a few other places. Is this correct?\\

\response{See response to SC-2 above.}


\noindent \ul{\textit{\textbf{SC-8:}}} On page 13 could you say a bit more about the experimental outcomes and theoretical guarantees presented in (2.3) and (2.4) in terms of theta and eta?\\

\response{Thank you for the suggestion.  For this experiment, $\theta$ (the angle between $Ax$ and $b$) is zero in exact arithmetic, which means we expect $\tan \theta \approx 0$ and the forward error of QR is governed by $\kappa \cdot \epsilon_{mach}$.  We have added a sentence to the discussion in Section 4.2 to make this point and better explain the numerical observations.  We also computed $\theta$ and $\eta$ for all the right hand sides of the experiment and observed that on average, $\theta$ values are around the square root of machine precision and $\eta$ values get quite large, up to $10^{14}$ for the most ill-conditioned problem.  Both of those drive down the magnitude of the second term in the error bound.}

\noindent \ul{\textit{\textbf{SC-9:}}} For the computational experiments, it is nice that a GitHub repository is included for reproducibility. It would be helpful if the authors could share a bit of details about the code in the text as well (e.g., it is written in MATLAB) as well as the machine used for testing (just to ensure the reader that the tests were comparing apples-to-apples).\\

\response{We added a description of the computational setup to the beginning of Section 4.}

\noindent \ul{\textit{\textbf{SC-10:}}} For the randomized example why was the convergence tolerance set to $10^{-15}$? Was this to ensure the experiments ran the maximum iteration number of 20?\\

\response{That is correct, the convergence tolerance detail is superfluous.  We re-worded to avoid confusion.}


\bibliographystyle{plain}
\bibliography{../reference}
\end{document}
